\section{Introduction}\label{sec:offline-intro}

Central-bank digital currency (CBDC) is a form of central bank money, on par
with cash and central-bank reserves. In CBDC, central bank money is digitally
represented as tokens, with each token encoding an owner and a
value. CBDCs are either wholesale or retail: the former is
a payment instrument accessible only to financial institutions to settle their
transactions, while the latter is available to the general public to conduct
retail payments.

With over 100 central banks currently exploring CBDC, and a
few launching pilot programs and developing regulatory frameworks, CBDC systems
may become a reality in the medium-term future~\cite{cbdctracker}. This interest is due in large
part to the pressing need to address inefficiencies prevalent in today's
financial markets and payment systems. Wholesale CBDC holds the potential to
mitigate settlement delays, reduce counter-party risks, and expedite
cross-border transactions. Retail CBDC is anticipated to slash transaction
fees, ignite innovation in digital payments, and foster financial inclusion.

Retail CBDC is the most challenging to realise given the stringent
requirements around resilience, scalability, privacy, and support for
offline payments. Offline payments are a pre-requisite to digital
financial inclusion, whose goal is to empower individuals to conduct digital
payments irrespective of their location or the quality of their network
coverage, and they are necessary for resilient and robust digital payments
that can operate even during blackouts or power outages. Offline payments
assume a setting where the payer and the payee are offline, except when they
withdraw or deposit their CBDCs.

To facilitate adoption against cash or other offline payment instruments, CBDC
offline payments must mimic some of the properties of physical cash. More
specifically, offline payments must be (1) fast: an offline payment
should take only a few seconds to complete; (2) transferable: CBDCs
received offline can be used in subsequent offline payments without
restriction; and (3) privacy-preserving: the central bank and the
commercial banks (the intermediaries) should not learn any information about
an offline payment beyond what is necessarily leaked at the withdrawal and
deposit of the CBDCs involved, such as the value of and identities of the
participants.

\paragraph{The challenge of offline double spending.} The complexity of
offline payments stems from the inability of the payee to verify in real time
whether the tokens they receive have already been spent. Without this
verification, double spending is trivial: a malicious payer can spend the same
token, offline, over and over again without restriction. Curtailing this
ability can be pursued through preventive or deterrent measures,
or a combination of both. A preventive measure is a mobile app that forbids
double spending by design, deleting a token as soon as it is sent; the caveat
is that an attacker can tamper with the app and bypass such restrictions. An
alternative is to leverage a secure element so that these restrictions
cannot easily be bypassed. A deterrent measure instead identifies double
spenders when double spending is detected at the time of deposit; combined
with the right punitive measures, this can discourage double spending by
individuals, though it will not deter state actors.

\paragraph{Our approach.} We combine secure elements and the identification of
double spenders to achieve the highest levels of security. The secure element
guarantees that double spending attacks are expensive and cannot be executed
at scale, while the identification of double spenders ensures that, on the
off-chance an attack succeeds, the responsible party is identified and held
accountable. A secure element, however, has limited storage and compute. While
it can accommodate a solution that uses ECDSA signatures to authenticate
CBDCs and authorise offline payments, such a solution fails to meet the
privacy requirements above: ECDSA signatures leak the identities of the payer
and the payee unless one-time public keys are used, and one-time public keys
in turn make double spender identification difficult, if not impossible.

We instead rely on (1) anonymous credentials to allow the payer and
the payee to transact anonymously, such that no one else can learn their
identities, (2) blind signatures, and (3) zero-knowledge
proofs of knowledge of signatures to ensure that neither the central bank nor
the intermediaries can link a withdrawal to a deposit. Finally, to
efficiently identify double spenders, we use (4) verifiable
encryption to encrypt the identity of the payer for an authorised auditor,
who is tasked with decrypting ciphertexts and de-anonymising double spenders.

These cryptographic primitives are too expensive for off-the-shelf secure
elements. We address this by splitting the computation of sending and
receiving an offline payment into two categories: one that is
computationally heavy but requires no secret keys, and one that is
lightweight, comparable to an ECDSA signature, and makes use of a
signing key. The heavy computation is outsourced to a mobile device, and the
lightweight computation is performed on the secure element, which -- in
addition to authorising payments with its signing key -- is also responsible
for ensuring that a token is used only once. To demonstrate the viability of
this approach, we benchmark it using smart cards as a stand-in for the secure
element. Even with 32 hops in a payment chain, both withdrawal and payment
complete comfortably within a second, and the secure element authorises a
payment in under $400\,\mathrm{ms}$.

\subsection{Related work}\label{sec:offline-related}

\paragraph{Original e-cash.} E-cash has been studied for several decades,
beginning with Chaum's solutions, which focus on the
privacy of users and on double spending detection~\cite{C:Chaum82,C:ChaFiaNao88}. Privacy is achieved through
RSA blind signatures, which break the link between withdrawals and deposits,
while double spending detection is realised using unique serial numbers.
Although efficient, these solutions guarantee neither transferability nor
theft prevention.

\paragraph{Divisible e-cash.} Divisible e-cash allows a user to withdraw a
token and split it among different recipients across multiple independent
transactions~\cite{C:OkaOht91,C:Okamoto95}. This comes at the cost of higher
complexity for withdrawals, payments, and deposits, which scale with the token
value. This is partially addressed by work that achieves constant computation
and communication complexity for withdrawal and payment, though deposits
still carry large communication and computational
costs~\cite{PKC:CPST15,ACNS:CPST15}. While divisibility is a useful feature, we leave it as
future work and focus instead on transferability.

\paragraph{Transferable e-cash.} Some schemes present e-cash that is both
transferable and divisible, at the cost of high computational complexity; we
instead focus on efficiency and scalability~\cite{C:OkaOht91,EC:ChaPed92}.
Fuchsbauer et al.\ maintain a constant token size regardless of
the number of transfers, but require user interaction to identify double
spenders, and when double spending is detected all previous owners of the
token also become identifiable~\cite{CANS:FucPoiVer09}. Baldimtsi et al.\
propose a scheme where token size grows linearly with the number of
transfers, but which reveals only the double spender's
identity~\cite{PKC:BCFK15}. Owing to the cost of this approach, we instead
rely on a trusted auditor, which can be distributed, to de-anonymise payments
as needed.

Some schemes introduce coin transparency, which ensures a payee
cannot tell whether they are receiving a token they previously
owned~\cite{ACNS:CanGou08,PKC:BauFucQia21}. One such scheme achieves this
through extensive use of zero-knowledge proofs together with a trusted party
that tracks every token in the system and can identify double spending
attempts~\cite{AFRICACRYPT:BCFGST11}. While this captures the strongest level of
privacy, we do not target coin transparency, owing to its computational
overhead and to the fact that physical cash does not satisfy it either.

The bounded-anonymity model keeps
transactions anonymous as long as each user interacts with a given merchant no
more than a fixed number of times; beyond this, the user becomes
identifiable~\cite{SCN:CamHohLys06}.
We instead let the conditions under which anonymity is revoked be defined
outside of the protocol itself.

\paragraph{Off-chain settlement.} Payment channels such as the Lightning
network and Teechain provide off-chain payments rather than offline
payments~\cite{poon2016bitcoin,lind2019teechain}. The Lightning network was
the first payment channel network designed for
Bitcoin: to improve latency and scalability, users
commit funds to an account controlled by signatures from both payer and
payee, then exchange small transactions off-chain, with only the final
batched transaction committed to the ledger~\cite{poon2016bitcoin}. A network of such channels lets
users route transactions through counterparties with whom they do not share
an account directly. Onion routing offers some notion of privacy, though it
has been shown that ordinary use of the network can bypass these privacy
aims~\cite{FC:KYPKDM21}. Teechain instead has users fund an account
controlled by a committee of treasurers, who use trusted hardware to store
their share of the private key and so raise the cost of stealing funds, in
order to use that value off-chain~\cite{lind2019teechain}. Teechain does not
target privacy, and is motivated by reducing the on-chain storage costs of
every transaction rather than by enabling offline payments.

\paragraph{CBDC offline payments.} Several reports motivate the need for CBDC
offline payments and describe current technologies. Project
Polaris\footnote{Project Polaris: secure and resilient CBDC systems, offline
and online.
\url{https://www.bis.org/about/bisih/topics/cbdc/polaris.htm}.} from the Bank
for International Settlements sets out design goals for offline payment
systems, distinguishing fully offline systems, where settlement occurs
offline; intermittently offline systems, where settlement is generally
offline but occasional synchronisation with the ledger is required once, for
example, the total value transferred crosses a risk threshold; and staged
offline systems, where value is exchanged offline but not settled until the
payee comes online.

Sweden has piloted its e-krona project on the Corda
platform, focusing on legal analysis
and integration with point-of-sale terminals; locally stored e-krona can be
used offline and, like cash, is not recoverable if the wallet is lost, and
transactions formed offline are not settled until the device goes online due
to online double-spending checks~\cite{brown2016corda,soderberg2019krona}. The People's Bank of China has designed a
system for digital payments\footnote{Reported at
\url{https://www.nytimes.com/2021/03/01/technology/china-national-digital-currency.html}.}
in which transaction values are hidden, though user anonymity and
transaction unlinkability do not appear to be implemented, and offline
payments are not specifically addressed. The Bank of Ghana's design document
for the digital Cedi (eCedi) highlights offline payments as a key focus
area\footnote{\url{https://www.bog.gov.gh/news/design-paper-of-the-digital-cedi-ecedi/}.}
while emphasising privacy alongside regulatory compliance. The US Federal
Reserve Board's report on offline payments
distinguishes hybrid payment systems, where the payee must come online to
settle, from fully offline systems, where this is not required, and finds no
production examples of the latter; it highlights the digital Rupee from the
National Payments Corporation of
India\footnote{\url{https://rbi.org.in/scripts/FAQView.aspx?Id=169}.}, which
allows offline transactions from a smartphone app with settlement once the
app reconnects, and Pix in Brazil, which allows the payer to be offline when
creating a transaction but requires the payee's point-of-sale terminal to be
online to authorise it~\cite{aboulaiz2024offline,lobo2021pix}.

\paragraph{Decentralised privacy-preserving payments.}
Peredi, Platypus, and
UTT give privacy-preserving token schemes that
distribute trust via threshold cryptography while providing accountability
and compliance guarantees, but do not address the offline
setting~\cite{CCS:KiaKohSar22,CCS:WKDC22,EPRINT:TBAAGP22}. OPERA
enables offline payments and prevents double spending using one-time readable
memory, but this technology is not widely available in
practice~\cite{park2017opera}. Rial and Piotrowska distribute the role of the bank via threshold
issuance and consider divisible offline payments, while we
focus on efficient transferability~\cite{PoPETS:RiaPio23}. PayOff offers a
transferable offline CBDC solution using more complex cryptography than ours,
and assigns the central bank responsibility for covering double-spent values,
potentially inflating the total value in the system; we instead prioritise
equipping users with a secure element to prevent double spending rather than
only detect it~\cite{Beer2024}.

\paragraph{Our contribution.} We give a transferable, privacy-preserving
offline CBDC payment scheme that splits the cost of authorising a payment
between a lightweight secure element and a more powerful mobile device, so
that the secure element need only perform operations comparable in cost to
producing an ordinary digital signature. Anonymous credentials, blind
signatures, and zero-knowledge proofs of knowledge hide the identities of
payers and payees from the central bank and from the intermediaries, while a
verifiable encryption under an auditor's key allows double spenders to be
efficiently identified without compromising the privacy of honest users. We
implement the scheme on commodity smart cards and mobile phones and show that
it meets the latency expected of everyday payments.

\paragraph{Layout.} Section~\ref{sec:offline-definitions} describes the
system's participants and trust assumptions and formalises the security goals
of an offline payment system as an ideal functionality.
Section~\ref{sec:construction} gives the construction, first at the level of
intuition and then in full detail. Section~\ref{sec:offline-instantiation}
instantiates the building blocks the construction relies on.
Section~\ref{sec:offline:proof} proves the construction secure and discusses
its security in practical deployments. Section~\ref{sec:eval} evaluates our
implementation, and Section~\ref{sec:offline-summary} concludes.
